\documentclass[11pt]{article}
\usepackage[margin=1in]{geometry}
\usepackage{amsmath,amssymb}
\usepackage{graphicx}
\usepackage{booktabs}
\usepackage{hyperref}
\usepackage{microtype}

\title{\textbf{Paper 37: Wave Radius Controls Zone Resolution Bandwidth;\\
K-Stage Relay Chains Do Not Compound Structure}}
\author{VCSM Theoretical Series}
\date{2026-03-10}

\begin{document}
\maketitle

\begin{abstract}
Paper~36 established that in the wave-dominated regime, the effective coupling
length is $\xi_\mathrm{eff} \approx r_\mathrm{wave}$, so the critical zone number
satisfies $N_\mathrm{crit} \approx L/(4r_\mathrm{wave})$.
We test this directly by sweeping $r_\mathrm{wave} \in \{1,2,3,4\}$ and measuring
relay gain $G$ at zone counts straddling the predicted threshold.
Experiment~A confirms $N_\mathrm{crit} \propto r_\mathrm{wave}^{-1}$ with an
empirical fit $N_\mathrm{crit} \approx L/(5r_\mathrm{wave})$; at $r_\mathrm{wave}=4$
the over-perturbation regime renders $G<1$ at all tested $N$.
Experiment~B tests whether a $k$-stage relay chain amplifies structure
geometrically; it does not. Gain $G_k$ remains approximately equal to the
single-stage gain $G_1 \approx 2$ for all $k \in \{1,\ldots,4\}$ and positional
noise $\sigma \in \{0,2,5\}$, confirming that the bimodal attractor saturates at
the first stage and subsequent stages merely inherit it.
\end{abstract}

\section{Introduction}

The relay coupling framework developed in Papers~35--36 establishes two key
dimensionless parameters governing geographic information transfer in VCML:
\begin{align}
  \Phi &= \frac{WR \cdot r_\mathrm{wave}}{T_\mathrm{dur}}, \quad
  \zeta = \frac{z_w}{\xi_\mathrm{eff}},
\end{align}
where $\Phi$ is the wave flux (waves per site per cycle) and $\zeta$ is the zone
resolution (zone width in units of the effective coupling length).
Paper~36 showed that $\xi_\mathrm{eff} \approx r_\mathrm{wave}$ in the
wave-dominated regime (rather than $\xi_\mathrm{diffusion} \propto \sqrt{D}$),
so the threshold becomes
\begin{equation}
  \zeta_\mathrm{crit} = z_w^* / r_\mathrm{wave} = \text{const},
  \qquad
  N_\mathrm{crit} = L / (z_w^* \cdot r_\mathrm{wave}^{-1}).
\end{equation}

This paper answers two residual questions:
\begin{enumerate}
  \item \textbf{(Q-rw)} Does $N_\mathrm{crit}$ scale as $r_\mathrm{wave}^{-1}$
        as predicted, with the same proportionality constant across all tested
        $r_\mathrm{wave}$ values?
  \item \textbf{(Q-chain)} Does coupling $k$ lattices in a chain amplify
        structure geometrically (i.e., does $G_k$ grow with $k$), or is gain
        saturated at the first stage?
\end{enumerate}

\section{Methods}

\subsection{Experiment A: Wave-Radius Sweep}

We sweep $r_\mathrm{wave} \in \{1,2,3,4\}$ at fixed $WR=4.8$, varying $N_\mathrm{zones}$
to straddle the predicted threshold $N_\mathrm{crit}^\mathrm{pred} = L/(4r_\mathrm{wave})
\in \{10, 5, 3.3, 2.5\}$.
Wave activation is parameterized as
\begin{equation}
  a(d) = \max\!\left(0,\; 1 - \frac{0.8 \cdot d}{r_\mathrm{wave}}\right), \quad d \le r_\mathrm{wave},
\end{equation}
so that the site at the edge of the footprint ($d = r_\mathrm{wave}$) always
receives activation $a = 0.2$, ensuring the footprint boundary is consistent
across $r_\mathrm{wave}$ values.
Relay gain $G = \mathrm{sg4n_{geo}} / \mathrm{sg4n_{ctrl}}$ is measured as
in Paper~35; $N_\mathrm{crit}$ is defined as the largest $N$ for which $G > 1.5$.
Five seeds per condition; all other parameters at VCML standard values.

\subsection{Experiment B: K-Stage Relay Chain}

We construct chains of $k \in \{1,2,3,4\}$ VCML lattices $L_1 \to \cdots \to L_k$,
where $L_1$ receives geographic waves (correct zone class) and each subsequent
stage $L_{i+1}$ receives a copy of $L_i$'s wave positions displaced by
Gaussian noise $\mathcal{N}(0, \sigma^2)$ independently per axis, with
$\sigma \in \{0, 2, 5\}$ (in lattice sites).
Zone class for the displaced position is recomputed from the new coordinates.
Standard parameters: $WR=4.8$, $N_\mathrm{zones}=4$, $r_\mathrm{wave}=2$.
Stage-$k$ relay gain is $G_k = \mathrm{sg4n}(L_k) / \mathrm{sg4n}(\mathrm{ctrl})$,
where ctrl is a single unlinked lattice receiving random-position waves.
Five seeds per condition; 275 total runs.

\section{Results}

\subsection{Experiment A: $N_\mathrm{crit} \propto r_\mathrm{wave}^{-1}$}

\begin{table}[h]
  \centering
  \caption{Relay gain $G$ by wave radius and zone count. Bold marks largest
           $N$ with $G > 1.5$ (i.e., $N_\mathrm{crit,obs}$).}
  \label{tab:rwave}
  \small
  \begin{tabular}{cccccccc}
    \toprule
    $r_w$ & $N$ & $z_w$ & $z_w/r_w$ & $G$ & $\pm$SE
          & $N_\mathrm{crit}^\mathrm{pred}$ & $N_\mathrm{crit}^\mathrm{obs}$ \\
    \midrule
    1 & 20 & 2 & 2.0 & 1.369 & 0.148 & & \\
    1 & 10 & 4 & 4.0 & 1.402 & 0.197 & & \\
    1 & \textbf{8} & 5 & 5.0 & 1.769 & 0.269 & 10 & \textbf{8} \\
    1 & 5  & 8 & 8.0 & 1.708 & 0.240 & & \\
    \midrule
    2 & 8 & 5 & 2.5 & 0.937 & 0.121 & & \\
    2 & 5 & 8 & 4.0 & 1.061 & 0.094 & & \\
    2 & \textbf{4} & 10 & 5.0 & 2.669 & 4.683 & 5 & \textbf{4} \\
    2 & 2 & 20 & 10.0 & 6.187 & 7.882 & & \\
    \midrule
    3 & 5 & 8 & 2.67 & 0.468 & 0.077 & & \\
    3 & 4 & 10 & 3.33 & 0.531 & 0.195 & & \\
    3 & \textbf{2} & 20 & 6.67 & 13.985 & 11.994 & 3.3 & \textbf{2} \\
    \midrule
    4 & 5 & 8 & 2.0 & 0.536 & 0.083 & & \\
    4 & 4 & 10 & 2.5 & 0.411 & 0.189 & 2.5 & --- \\
    4 & 2 & 20 & 5.0 & 0.339 & 0.151 & & \\
    \bottomrule
  \end{tabular}
\end{table}

Table~\ref{tab:rwave} and Figure~1a show that $G$ drops sharply below 1.5
at approximately the same threshold in $z_w/r_w$ space across all three
resolved values of $r_\mathrm{wave}$.
The observed critical zone counts are $N_\mathrm{crit,obs} \in \{8, 4, 2\}$
for $r_\mathrm{wave} \in \{1, 2, 3\}$, confirming the $r_\mathrm{wave}^{-1}$
scaling with ratio $8:4:2$ matching $1:2:3$ exactly.

\paragraph{Empirical threshold.}
The original Paper~35 prediction was $N_\mathrm{crit} = L/(4r_\mathrm{wave})$
(i.e., $z_w > 4\xi_\mathrm{eff} = 4r_\mathrm{wave}$).
The observations are systematically 20\% lower:
$N_\mathrm{crit,obs}/N_\mathrm{crit,pred}^{(4)} = 0.80$ for $r_\mathrm{wave} \in \{1,2\}$.
The corrected empirical fit is
\begin{equation}
  \boxed{N_\mathrm{crit} \approx \frac{L}{5\, r_\mathrm{wave}}},
  \qquad \text{i.e., threshold at } z_w / r_\mathrm{wave} \approx 5.
\end{equation}
This corresponds to $\zeta_\mathrm{crit} = z_w^*/\xi_\mathrm{eff} = 5$
rather than 4.
Geometrically, a zone must be at least 5 wave-radii wide for geographic
relay to benefit structure. The slight upward revision from $4\to 5$ may
reflect a boundary-condition correction: sites near the zone edge receive
activation from both the zone's own wave and (weakly) from adjacent-zone waves,
effectively widening the cross-contamination region beyond the nominal footprint.
Figure~1b shows the N_crit vs r_wave scaling with both predictions.

\paragraph{Over-perturbation at $r_w = 4$.}
At $r_\mathrm{wave}=4$, $G < 1$ for all tested $N$ values, including $N=2$
($z_w/r_w = 5$, which is at the threshold for smaller $r_w$).
This is not a zone-resolution failure --- it is over-perturbation.
The activation profile $a(d) = 1 - 0.2d$ integrates to a much larger total
activation per wave than smaller radii: the wave footprint covers a disk of
radius 4 with a cone profile, delivering $\sim\!10\times$ more total energy
per launch than $r_w=1$.
With $WR=4.8$ fixed, $r_w=4$ places the system in the $\Phi \gg 1$ regime
(wave flux $\Phi = WR \cdot r_w / T_\mathrm{dur} = 4.8\times 4/15 = 1.28$,
vs.\ $\Phi=0.64$ at standard $r_w=2$), where Paper~35 showed $G < 1$.
Thus, $r_w=4$ is over-perturbed even at the most favorable zone count.

\subsection{Experiment B: Chain Relay Does Not Compound}

\begin{table}[h]
  \centering
  \caption{Stage-$k$ relay gain $G_k$ and ratio $\mathrm{sg4n}(L_k)/\mathrm{sg4n}(L_1)$.
           Ctrl baseline: $\mathrm{sg4n_{ctrl}} = 0.1241$.
           $L_1$ gain: $G_1 = 2.088$ (same row in all sigma conditions).}
  \label{tab:chain}
  \small
  \begin{tabular}{ccccc}
    \toprule
    $\sigma$ & $k$ & $G_k$ & $\mathrm{sg4n}(L_k)/\mathrm{sg4n}(L_1)$ \\
    \midrule
    0 & 1 & 2.088 & 1.000 \\
    0 & 2 & 1.783 & 0.854 \\
    0 & 3 & 1.398 & 0.670 \\
    0 & 4 & 2.250 & 1.077 \\
    \midrule
    2 & 1 & 2.088 & 1.000 \\
    2 & 2 & 2.272 & 1.088 \\
    2 & 3 & 2.008 & 0.962 \\
    2 & 4 & 2.326 & 1.114 \\
    \midrule
    5 & 1 & 2.088 & 1.000 \\
    5 & 2 & 2.293 & 1.098 \\
    5 & 3 & 5.418 & 2.595 \\
    5 & 4 & 2.542 & 1.217 \\
    \bottomrule
  \end{tabular}
\end{table}

Table~\ref{tab:chain} shows that $G_k \approx G_1 \approx 2$ throughout for
$\sigma \in \{0, 2\}$, with no systematic increase or decrease with $k$.
The prediction that $G_k$ would \emph{compound} (i.e., $G_k > G_1$) is falsified.
The prediction that $\sigma > 0$ would cause \emph{geometric decay} (i.e.,
$G_k \sim G_1 (1 - \sigma/z_w)^{k-1}$) is also not confirmed --- $G_k$ at
$\sigma=2$ and $\sigma=5$ is statistically indistinguishable from $G_k$ at $\sigma=0$.

The $k=3$, $\sigma=5$ outlier ($G_3 = 5.42$) is an artifact of the bimodal
attractor combined with $n=5$ seeds: one seed in which $L_3$ happened to be
in the high-structure basin while $L_1$ was not inflates the mean.
The within-seed ratio at that condition includes an extreme observation
(ratio = 2.595), consistent with random basin-entry rather than systematic amplification.

\paragraph{Mechanism: saturation at the first stage.}
The bimodal attractor of VCML has two basins: high-structure ($G \gg 1$) and
low-structure ($G \approx 1$).
Geographic relay at stage 1 biases $L_2$ toward the high-structure basin.
Once there, $L_2$'s own copy-forward loop maintains structure independently
of further relay; \emph{additional stages relay the same already-structured
signal} and cannot increase basin occupancy beyond what stage 1 achieves.
This is a ceiling effect, not a degradation: geographic coupling is useful for
initial basin selection, not for sustained amplification.

\paragraph{Robustness to positional noise.}
Positional noise $\sigma \in \{2, 5\}$ (20--50\% of zone width $z_w = 10$)
does not measurably degrade relay effectiveness.
This confirms that the zone discrimination is robust: a wave landing within
the correct zone but $\pm\sigma$ sites from the ideal position still provides
a consistent zone-class signal, because the footprint radius ($r_w = 2$) is
much smaller than the zone width and the system averages over many waves per
epoch.

\section{Discussion}

\subsection{Updated N-Crit Formula}

Combining Papers~35--37, the zone-resolution constraint is:
\begin{equation}
  N_\mathrm{crit} = \frac{L}{5\, r_\mathrm{wave}},
  \qquad \text{subject to } \Phi = \frac{WR \cdot r_\mathrm{wave}}{T_\mathrm{dur}} < 1.
\end{equation}
The factor 5 (vs.\ the theoretical 4) is an empirical correction accounting
for zone-boundary wave contamination.
Both constraints must be satisfied simultaneously: the zone must be wide enough
for discrimination (factor-5 condition) and the wave flux must be below 1
to avoid over-perturbation.
For standard parameters ($L=40$, $r_w=2$, $T_\mathrm{dur}=15$, $WR=4.8$):
$N_\mathrm{crit} = 4$ and $\Phi = 0.64$ (below 1), consistent with all prior results.

\subsection{Limits of Chain Relay}

The null result for $G$-compounding has a positive implication: a two-stage
relay $(L_1 \to L_2)$ achieves the same structure quality as a ten-stage relay.
Geographic coupling is efficient (one hop suffices) and robust (positional
noise within zone width is tolerable).
The ceiling effect suggests that the bottleneck is not signal transmission
fidelity but basin-entry probability: the bimodal attractor has a fixed
probability of being in the high-structure basin under geographic coupling,
and additional reinforcement stages do not change that probability.

\subsection{Open Questions}

\begin{enumerate}
  \item \textbf{WR-dependent threshold}: at each $r_w$, there should be an
        optimal WR satisfying $\Phi = WR \cdot r_w / T_\mathrm{dur} < 1$.
        Does the peak gain $G_\mathrm{max}$ depend on $r_w$?
  \item \textbf{Direct basin measurement}: what fraction of seeds is in the
        high-structure basin under geo vs.\ ctrl coupling?
        This would directly measure the basin-entry effect and resolve the
        large SEs at $N=2$--4 (bimodal variance).
  \item \textbf{2D relay surface}: map $G(WR, N, r_w)$ jointly for the
        full parameter space.
\end{enumerate}

\section{Conclusion}

We confirm that $N_\mathrm{crit} \propto r_\mathrm{wave}^{-1}$ with empirical
constant $L/5$ (vs. the theoretical $L/4$), establishing $z_w > 5r_\mathrm{wave}$
as the zone-resolution condition for effective geographic relay.
The over-perturbation ceiling ($\Phi < 1$) is a second independent constraint
that determines the maximum usable $r_\mathrm{wave}$ at fixed $WR$.
Chain relay does not compound: $G_k \approx G_1$ for all tested $k$ and $\sigma$,
because geographic coupling selects the attractor basin at stage 1 and
subsequent stages merely inherit it.
These results close the quantitative theory of geographic relay coupling:
$N_\mathrm{crit} = L/(5r_w)$, $\Phi < 1$, one relay hop suffices.

\end{document}
